\documentclass[header={Lecture 6 Exam Study Notes}]{study-note}

\begin{document}

\StudyTitleBlock{06}{Lecture 6: Exam Study Notes}{Uniform boundedness, open mapping, bounded inverse, and closed graph}

\vspace{0.8em}

\begin{focusbox}
This lecture is one of the central exam blocks. These are the ``great theorems'' built on Baire's theorem and used everywhere later.
\end{focusbox}

\tableofcontents

\section{Uniform Boundedness Principle}

\begin{thmbox}{Theorem 8.5: Banach--Steinhaus uniform boundedness principle}
\thmBanachSteinhaus
\end{thmbox}

\begin{warningbox}{Common trap}
Banach--Steinhaus requires the domain space $X$ to be Banach. Pointwise boundedness alone is not enough on incomplete spaces.
\end{warningbox}

\begin{prooflevelbox}{SKETCH}
Know the structure: pointwise bounded $\Rightarrow$ there exists a ball on which the family is uniformly bounded (Baire category) $\Rightarrow$ rescale to get a global bound.
\end{prooflevelbox}

\paragraph{Meaning of the hypothesis.}
``Pointwise bounded'' means: fix one $x\in X$ and look at all values $Lx$ with $L\in\mathcal F$; those values stay uniformly bounded. Baire's theorem upgrades this pointwise information to a uniform operator-norm bound.

\paragraph{Main consequences.}
\begin{itemize}
\item weakly convergent and weak-* convergent sequences are bounded;
\item pointwise bounded families of operators are uniformly bounded;
\item the pointwise limit of bounded linear operators is bounded (next theorem).
\end{itemize}

\begin{thmbox}{Theorem 8.6: Pointwise limit of bounded linear operators is bounded}
\thmLimitOperatorBounded
\end{thmbox}

\begin{prooflevelbox}{SKETCH}
Apply Banach--Steinhaus to the sequence $(L_n)$, then verify the limit operator is linear and bounded.
\end{prooflevelbox}

\section{Open Mapping And Bounded Inverse}

\begin{thmbox}{Theorem 8.7: Open mapping theorem}
\thmOpenMapping
\end{thmbox}

\begin{prooflevelbox}{SKETCH}
Know that Baire category gives $B_Y(0,c)\subset \ol{L(B_X(0,1))}$ for some $c>0$, then iterate and scale to show the image contains an open ball.
\end{prooflevelbox}

\begin{exbox}{Quantitative form}
The open mapping theorem is often used through the equivalent statement:
there exists $c>0$ such that
\[
B_Y(0,c)\subset L(B_X(0,1)).
\]
After scaling, this gives a uniform estimate for solving $Lx=y$.
\end{exbox}

\begin{thmbox}{Theorem 8.8: Bounded inverse theorem}
\thmBoundedInverse
\end{thmbox}

\begin{prooflevelbox}{STATEMENT}
Direct consequence of the open mapping theorem applied to a bijective $L$.
\end{prooflevelbox}

\begin{thmbox}{Theorem 8.9: Sufficient condition for equivalence of norms}
\thmComparableNorms
\end{thmbox}

\begin{prooflevelbox}{STATEMENT}
Apply the bounded inverse theorem to the identity map.
\end{prooflevelbox}

\paragraph{Reason.}
Apply the bounded inverse theorem to the identity map between the two normed versions of $X$.

\section{Closed Graph Theorem}

\begin{defbox}{Closed linear operator}
For a linear operator $L:X\to Y$, its graph is
\[
\operatorname{Graph}(L)=\{(x,Lx):x\in X\}\subset X\times Y.
\]
The operator is \textbf{closed} if its graph is closed.
\end{defbox}

\begin{thmbox}{Theorem 8.10: Closed graph theorem}
\thmClosedGraph
\end{thmbox}

\begin{warningbox}{Common trap}
The closed graph theorem requires the operator to be defined on all of $X$. For partially defined operators, use the graph norm formulation instead.
\end{warningbox}

\begin{prooflevelbox}{SKETCH}
Know how to define the graph norm and use completeness of the graph to apply the open mapping theorem to the projection onto the first coordinate.
\end{prooflevelbox}

\paragraph{How to check closedness.}
To prove that $L$ is closed, one shows:
whenever $x_n\to x$ in $X$ and $Lx_n\to y$ in $Y$, then $y=Lx$.

\begin{thmbox}{Theorem 8.11: Graph norm characterization of closed operators}
Let $L:D(L)\subset X\to Y$ be linear, where $X,Y$ are Banach spaces. Define the graph norm on $D(L)$ by
\[
\norm{x}_{L}:=\norm{x}_X+\norm{Lx}_Y.
\]
Then $L$ is closed if and only if $D(L)$ is complete with respect to the graph norm.
\end{thmbox}

\paragraph{Why this matters.}
This is the right framework for operators that are not defined on all of $X$. The closed graph theorem is the special case $D(L)=X$.

\section{Hellinger--Toeplitz Connection}

\begin{warningbox}{Exam-specific use}
One oral exam topic asks for the use of the closed graph theorem in the proof of the Hellinger--Toeplitz theorem: an everywhere-defined symmetric operator on a Hilbert space is bounded.
\end{warningbox}

\section{Exam-Style Exercises}

\begin{exbox}{Exercise 6.1: Weakly convergent sequences are bounded}
Let $X$ be a Banach space and suppose $x_n\rightharpoonup x$ in $X$. Prove that $(x_n)$ is bounded in norm.

\paragraph{Solution pattern.}
For each $n$, define $T_n:X^*\to\K$ by $T_n(\varphi)=\varphi(x_n)$. For fixed $\varphi$, the scalar sequence $T_n(\varphi)$ converges, hence is bounded. Apply Banach--Steinhaus to $(T_n)\subset (X^*)^*$, then identify $\norm{T_n}$ with $\norm{x_n}$ through the canonical embedding.
\end{exbox}

\begin{exbox}{Exercise 6.2: Pointwise limit of bounded operators}
Let $X$ be Banach, $Y$ normed, and let $(T_n)\subset\mathcal L(X,Y)$. Assume that for every $x\in X$, the limit
\[
Tx:=\lim_{n\to\infty}T_nx
\]
exists in $Y$. Prove that $T\in\mathcal L(X,Y)$.

\paragraph{Solution pattern.}
Linearity follows by passing to the limit. Use Banach--Steinhaus to get $\sup_n\norm{T_n}<\infty$, then estimate
\[
\norm{Tx}_Y\le \sup_n\norm{T_n}\,\norm{x}_X.
\]
\end{exbox}

\begin{exbox}{Exercise 6.3: Bounded inverse theorem in action}
Let $X,Y$ be Banach spaces and let $T\in\mathcal L(X,Y)$ be bijective. Prove that there exists $C>0$ such that
\[
\norm{x}_X\le C\norm{Tx}_Y
\qquad \forall x\in X.
\]

\paragraph{Solution pattern.}
Apply the bounded inverse theorem to $T^{-1}$ and write
\[
\norm{x}_X=\norm{T^{-1}Tx}_X\le \norm{T^{-1}}\norm{Tx}_Y.
\]
\end{exbox}

\begin{exbox}{Exercise 6.4: Equivalence of two comparable complete norms}
Let $\norm{\cdot}_1$ and $\norm{\cdot}_2$ be norms on the same vector space $X$. Assume both normed spaces are Banach and
\[
\norm{x}_2\le C\norm{x}_1 \qquad \forall x\in X.
\]
Prove that the norms are equivalent.

\paragraph{Solution pattern.}
Consider the identity map from $(X,\norm{\cdot}_1)$ to $(X,\norm{\cdot}_2)$. It is bounded and bijective, so the bounded inverse theorem makes the inverse bounded as well.
\end{exbox}

\begin{exbox}{Exercise 6.5: Closed graph criterion}
Let $X,Y$ be Banach spaces and let $T:X\to Y$ be linear. Suppose that whenever $x_n\to x$ in $X$ and $Tx_n\to y$ in $Y$, then $y=Tx$. Prove that $T$ is bounded.

\paragraph{Solution pattern.}
The assumption says exactly that the graph of $T$ is closed in $X\times Y$. Then apply the closed graph theorem.
\end{exbox}

\begin{exbox}{Exercise 6.6: Every bounded operator is closed}
Let $X,Y$ be normed spaces and let $T\in\mathcal L(X,Y)$. Prove that $T$ is closed.

\paragraph{Solution pattern.}
If $(x_n,Tx_n)\to (x,y)$ in $X\times Y$, then continuity of $T$ gives $Tx_n\to Tx$. Uniqueness of limits in $Y$ implies $y=Tx$.
\end{exbox}

\begin{exbox}{Exercise 6.7: Why the full-domain assumption matters}
Give an example of a bounded linear map on a non-closed subspace whose graph is not closed as a subset of the ambient product.

\paragraph{Solution pattern.}
Take $X=Y=\R$, domain $D$ equal to the rational numbers, and define $Tq=q$. The map is bounded on $D$, but its graph
\[
\{(q,q): q \text{ rational}\}
\]
is not closed in $\R^2$.
\end{exbox}

\section{What To Memorize Precisely}

\begin{focusbox}
You should know the exact statements and standard consequences of:
\begin{enumerate}
\item Theorems \textbf{8.5}, \textbf{8.6}, \textbf{8.7}, \textbf{8.8}, \textbf{8.9}, \textbf{8.10}, \textbf{8.11};
\item boundedness of weakly convergent sequences from Banach--Steinhaus;
\item the quantitative ball-image form of the open mapping theorem;
\item continuity of inverse from the open mapping theorem;
\item equivalence of norms via the identity map;
\item the graph norm criterion for closed operators.
\end{enumerate}
\end{focusbox}

\end{document}
